% Λύση του 67ου προβλήματος της ενότητας «Άλυτα Προβλήματα».
\subsection*{Λύση Προβλήματος 67 --- Αποσβεσμένος ταλαντωτής}

\begin{alist}
  \item Για κάθε \(t\geq0\) ισχύει
        \[
          -1\leq \syn(\omega t)\leq1.
        \]
        Επειδή \(A>0\) και \(e^{-\gamma t}>0\), έχουμε
        \[
          -Ae^{-\gamma t}
          \leq
          Ae^{-\gamma t}\syn(\omega t)
          \leq
          Ae^{-\gamma t}.
        \]
        Δηλαδή
        \[
          -Ae^{-\gamma t}\leq x(t)\leq Ae^{-\gamma t}.
        \]

        Επειδή \(\gamma>0\),
        \[
          \lim_{t\to+\infty}Ae^{-\gamma t}=0
        \]
        και
        \[
          \lim_{t\to+\infty}\bigl(-Ae^{-\gamma t}\bigr)=0.
        \]
        Άρα, από το κριτήριο παρεμβολής,
        \[
          {\lim_{t\to+\infty}x(t)=0}.
        \]

        Φυσικά, αυτό σημαίνει ότι η απομάκρυνση του ταλαντωτή από τη
        θέση ισορροπίας τείνει στο μηδέν. Η ταλάντωση σβήνει.

  \item Έχουμε
        \[
          t e^{-\gamma t}=\frac{t}{e^{\gamma t}}.
        \]
        Όταν \(t\to+\infty\), προκύπτει μορφή
        \[
          \frac{+\infty}{+\infty}.
        \]
        Εφαρμόζουμε τον κανόνα de l'H\^opital:
        \[
          \lim_{t\to+\infty}\frac{t}{e^{\gamma t}}
          =
          \lim_{t\to+\infty}\frac{1}{\gamma e^{\gamma t}}.
        \]
        Επειδή \(\gamma>0\), έχουμε
        \[
          \lim_{t\to+\infty}\gamma e^{\gamma t}=+\infty.
        \]
        Άρα
        \[
          {
          \lim_{t\to+\infty}t e^{-\gamma t}=0
          }.
        \]

  \item Από το πρώτο ερώτημα βρήκαμε ότι
        \[
          \lim_{t\to+\infty}x(t)=0.
        \]
        Επομένως, από τον ορισμό της οριζόντιας ασύμπτωτης, η ευθεία
        \[
          {y=0}
        \]
        είναι οριζόντια ασύμπτωτη της γραφικής παράστασης της \(x\)
        όταν \(t\to+\infty\).
\end{alist}

\paragraph{Πηγές και έλεγχος ρεαλιστικότητας}
Το μοντέλο αποσβεννυμένης ταλάντωσης είναι πραγματικό μοντέλο φυσικής:
το \en{OpenStax College Physics} εξηγεί ότι σε πραγματικούς ταλαντωτές
η τριβή και άλλες μη συντηρητικές δυνάμεις αφαιρούν ενέργεια και το
πλάτος μειώνεται με τον χρόνο \cite{openstaxCollegePhysicsDampedMotion}.
Ο παράγοντας \(e^{-\gamma t}\) εκφράζει ακριβώς αυτή τη μείωση του
πλάτους. Επειδή η τριγωνομετρική συνάρτηση παραμένει ανάμεσα στο
\(-1\) και στο \(1\), ενώ η περιβάλλουσα \(Ae^{-\gamma t}\) τείνει στο
μηδέν, το συμπέρασμα \(x(t)\to0\) είναι ρεαλιστικό: ο ταλαντωτής
τελικά πλησιάζει τη θέση ισορροπίας.
